\par\SourceLeaf{140}
Since this opinion and the interpretation given to it form the actual foundation of the arrangement which is the main subject of the 1st division of my memorandum, namely the rule of progression, I shall subject it further below to a detailed and exact examination, and therefore confine myself here, for the time being, to the following remarks.

The records do not show why \textsc{Kritter} speaks at the beginning of his opinion of 22—24 married professors, when P. had expressly mentioned 26, and only that number, in his letter. I suspect, however, that \textsc{Kritter} took those numbers 22—24 as applying to an \textit{earlier} time. To my knowledge, however, no complete registers of the personal circumstances of the members of the widows' fund have been kept in a form from which the number of \textit{married members} at any given time could be obtained. Thus either \textsc{Kritter} merely \textit{inferred} those numbers by a rough estimate from the number of all participants at an earlier time, or they rest on particular communications, which by the nature of the matter can then refer only to times not many years before, still remembered.

\textsc{Kritter}'s answer to the first question thus amounts, briefly, to saying that for the \textit{time} at which the maximum had occurred one could reckon one widow for about two existing marriages, hence for those 22—24 marriages 12 widows; this, he states, had been confirmed by the average of the last 17 years insofar as there had sometimes been more and sometimes fewer than 12 widows. That the number of widows corresponding to the present number of 26 marriages would then be one greater is not expressly stated, but is implicitly contained in this. In answer to the second question he states that, according to the experience of the Bremen widows' fund, where orphans' pensions last only until their \(18^{\text{th}}\) year, one reckons on an increase in the number of pensions by a sixth part; for the fund here, therefore, where their duration is to be 2 years longer, one would presumably have to assume somewhat more. — For the 24 marriages we thus arrive at somewhat more than 14, and for the 26 marriages, once their effect has been fully felt, at 15 according to \textsc{Kritter}'s words, or at somewhat more than \(15\tfrac12\), that is, nearly 16, according to the principles he has stated.

Whether, \textit{quite apart from a closer examination of the content of the opinion}, such a treatment of the subject can be called a \textit{detailed} and \textit{exact} investigation, as \textsc{Kaestner} had considered necessary, I leave aside here. On the basis of this opinion, however, P. now drafted a P. M. in which he shows that if the increase in pensions proposed by the governing board were fixed uniformly for all at 130 Rth., without a further special increase for the six oldest widows, this could be done without any danger to the fund even if the annual contributions were not increased; but that, if the increase in contributions in the manner indicated by the governing board were adopted, the extension of the orphans' pensions could also be introduced with all the greater safety, because in the prevailing circumstances this alteration could not be expected to take effect soon. P. finally closes his submission, to which I have no hesitation in giving my full assent \textit{up to this point}, with the following short addition. Since the subject of my own investigation, namely the rule of progression, here appears on the scene \textit{for the first time}, I copy it here completely and faithfully in P.'s own words:

\begin{quote}
‘In a general increase of the pensions to 130 Rth. there seems to me to be not the slightest danger:
\end{quote}
\SourceLeaf{141}
\begin{quote}
Rather, it also seems to me
\begin{list}{}{\setlength{\leftmargin}{1.8em}\setlength{\labelwidth}{1.4em}\setlength{\labelsep}{.4em}}
\item[3)] possible and useful that it should now be permissible to establish as a permanent rule that the widows' pensions are to be increased each time by 10 Thl. whenever the capital fund has grown by 5000 Thl. and the number of widows has not yet risen above the maximum of 15. It is clear that one can do this, for an increase of 10 Thl. for 15 widows amounts to 150 Thl., and 5000 Thl. at 3 percent yield just as much interest. But that the next time for an increase may come soon, even if our fund receives no extraordinary inflows, can at least be calculated as highly probable. Since at present we have only 10 widows to pension, over 1000 Thl. must remain in the fund every year if the new addition to the contributions is approved, and consequently 5000 Thl. must already have been added to the capital fund in 5 years, if the number of widows has not increased in the meantime; but even if one assumes the highly improbable case that the number were to increase each year by one widow until it had reached the maximum of 15, it could still scarcely take as long as 10 years.
\end{list}
It might at least be good if this circumstance were mentioned in the report to the Royal Government.’
\end{quote}

\noindent\hspace*{\parindent}Since what is necessary concerning the essential content of this article will be given further below in examining \textsc{Kritter}'s opinion and the conclusions attached to it, and in other places, only a few incidental circumstances will be touched upon here.

I. What is striking, coming from the pen of the curator of the widows' fund, is the \textit{incorrect} statement of the number of widows, or rather of pensions. At that time there were not 10 widows, nor 10 pensions, but 8 widows and, counting the children of Professor B., who had died on 8 June 1794, a total of 9 pensions. Nor can this inaccuracy be explained by supposing that the curator had not yet learned of the loss of the widow (M.) who had most recently died away from here (at Halle); for it is found that the amount of the pension for the last half year (Michaelmas 1793 to Easter 1794) was collected on the heiress's receipt, countersigned by P. and dated 3 May 1794.

II. The closing line (‘It might at least be good, etc.’) leaves us somewhat in the dark as to what the copied article is actually to be taken for: a formal motion put forward for a decision, or merely an idea thrown out. Neither the rescript nor the preceding deliberations had spoken of a definite rule for the variable pension rate applying for all future time. This was therefore a \textit{fourth} subject, newly added to those already under deliberation, and one which, being in no way less important than the other three, required the most careful examination from every side. Anyone putting forward a motion directed specifically to such an end is surely aware of the importance of the matter, which is scarcely brought out sufficiently by describing so vital a question as ‘this circumstance’. In the case of a remark, however, which has merely the character of a thought thrown out incidentally, one is more indulgent towards an inaccuracy that has slipped through inattention, and towards wording that is still defective. The numbering (3) should not mislead us here. In his paper P. does not \textit{count} the proposals, but the conclusions he derives from the rough calculations given beforehand. His No. 1 contains \textit{the} conclusion that all objections raised or to be raised against extending the orphans' pensions are disposed of if the proposed increase in contributions is adopted; and his No. 2, that the general increase of the pensions to 130 Rth. could indeed appropriately take place at once, whereas the exceptional increase to 150 Rth. for the six oldest widows would be neither just nor advisable.

